\scanpage{262}\mspage{\unsure{178}}
% archive sheet number: 261
que le foncteur canonique
\[
\operatorname{Ind}_{\pi}(\mathcal F^\pi_{\deleted{\unsure{\pi}}})\hookrightarrow M
\tag*{(4.7.11.1)}
\]
prolongeant l'inclusion $\mathcal F^\pi\subset M$, de façon
à commuter aux $\varinjlim$ \unsure{grdes} devant $\pi$, est
\uline{pleinement} \uline{fidèle}. De plus, lorsque
$\pi^{\pi_0}=\pi$, la condition c) des \unsure{fcg|fag} nous dit
\inserted{entre autres} que ce foncteur est ess. surjectif, donc
que c'est une \uline{équivalence de catégories} :
\[
\operatorname{Ind}_{\pi}(\mathcal F^\pi)\xrightarrow{\approx}M_{\unsure{\sim|a}}\quad
(\text{si }\pi=\pi^{\pi_0}).
\tag*{(4.7.11.2)}
\]

\uline{Proposition} 4.7.12. Soit $(\mathcal F^\pi)_{\pi\geq\pi_0}$ une fcg
\inserted{accessible (\unsure{fcga})} sur la $\unsure{U|M}$-cat. $M$, et soit $\pi$ un cardinal $\geq\pi_0$.
Les conditions suivantes sont équivalentes

a) Le foncteur pl. fidèle (4.7.11.1) est
ess. surjectif, i.e. est une équivalence de
catégories.

b) Pour \unsure{tout} $F\in\operatorname{Ob}M$, $\mathcal F^\pi_{/F}$ est
une catégorie \inserted{filtrante et} grande devant $\pi$.

\deleted{\unsure{On}} Quand ces conditions sont satisfaites, \inserted{et si $\pi>\pi_0$} on a
\deleted{Résulte aussitôt de 4.4.4.}
\[
\mathcal F^\pi=M_\pi\mathpunct{.}
\tag*{(4.7.12.1)}
\]

\uline{Dém.} L'équivalence de a) et b) résulte
aussitôt de 4.4.\unsure{2|8}. et de la définition
des catégories \unsure{grdes} devant $\pi$ (p. 81) : Que ces
